\begin{lemma}
Let $G$ be an $(n,d,\lambda)$-graph, allowing loops, and assume that its
adjacency operator has nontrivial spectral norm at most $\lambda$.  Then for
all vertex subsets $A,B$,
\[
 \left|e_G(A,B)-\frac d n|A||B|\right|
 \leq \lambda\sqrt{|A||B|}.
\]
The edge count is ordered, so a loop is counted when its endpoint lies in
both sets.  In the Lean interface, the spectral-norm hypothesis is written
as the displayed bilinear inequality for arbitrary real vectors, and the
conclusion is obtained by substituting the indicator vectors of $A$ and
$B$.
\end{lemma}
\begin{proof}
Let $1_A$ and $1_B$ be the two indicator vectors.  Substitute them into the
bilinear spectral inequality in the statement.  The bilinear term is exactly
the number of ordered adjacent pairs from $A$ to $B$, namely $e_G(A,B)$.
The two linear sums are $|A|$ and $|B|$, and the two squared Euclidean norms
are $|A|$ and $|B|$.  The resulting inequality is precisely the displayed
expander-mixing estimate.  No step uses irreflexivity, so the same argument
also covers allowed loops.
\end{proof}
